\section{Sintesis Literatur dan Posisi Penelitian}

Berdasarkan tinjauan pustaka yang telah dibahas, ABAC merupakan model kontrol akses yang relevan untuk lingkungan IoT karena mampu mendukung keputusan akses yang bersifat \textit{fine-grained} dan \textit{context-aware}. ABAC memungkinkan keputusan akses tidak hanya ditentukan berdasarkan identitas atau peran, tetapi juga berdasarkan atribut subjek, objek, aksi, serta kondisi lingkungan. Namun, penerapan ABAC pada IoT masih menghadapi tantangan karena evaluasi kebijakan membutuhkan atribut yang akurat dan mutakhir. Semakin banyak atribut yang digunakan, semakin besar pula biaya pengambilan atribut, komunikasi dengan PIP, dan latensi keputusan akses.

Pemindahan evaluasi keputusan akses ke \textit{edge} dapat mengurangi ketergantungan terhadap \textit{cloud} dan mempercepat proses otorisasi. Akan tetapi, pendekatan ini juga menimbulkan persoalan baru, terutama ketika atribut disimpan secara lokal pada \textit{cache gateway}. Atribut yang tersimpan di \textit{cache} dapat menjadi tidak mutakhir ketika terjadi perubahan kondisi perangkat, perilaku akses, status jaringan, \textit{firmware, role}, atau pencabutan hak akses. Kondisi tersebut dapat menyebabkan keputusan otorisasi yang tidak akurat, seperti \textit{false allow, false deny}, dan \textit{revocation false allow}.

Penelitian terdahulu telah membahas ABAC untuk IoT, otorisasi berbasis \textit{edge/fog}, sinkronisasi atribut, \textit{caching}, dan optimasi metaheuristik pada sistem IoT. Namun, masih terdapat celah penelitian pada optimasi \textit{freshness} atribut dalam ABAC berbasis \textit{edge}, khususnya dalam penentuan atribut yang boleh di-\textit{cache}, nilai TTL per atribut, dan ukuran \textit{cache} lokal secara terpadu dengan mempertimbangkan \textit{trade-off} antara performa dan kualitas keputusan otorisasi.

Posisi penelitian ini berada pada irisan antara ABAC untuk IoT, otorisasi berbasis \textit{edge}, manajemen \textit{freshness} atribut, dan optimisasi berbasis metaheuristik. Berbeda dari pendekatan yang hanya berfokus pada pemindahan evaluasi ke \textit{edge} atau penggunaan \textit{cache} untuk menurunkan latensi, penelitian ini menempatkan \textit{cache} atribut sebagai komponen keamanan yang dapat memengaruhi kualitas keputusan akses. Oleh karena itu, optimasi tidak hanya diarahkan untuk mempercepat evaluasi, tetapi juga untuk menekan risiko penggunaan atribut yang sudah usang.

Kebaruan penelitian ini terletak pada formulasi optimasi \textit{freshness} atribut pada ABAC berbasis \textit{edge} dengan pendekatan metaheuristik. Optimasi dilakukan terhadap \textit{cacheable attribute}, TTL per atribut, dan ukuran \textit{cache}, tanpa mengubah semantik kebijakan ABAC. Dengan demikian, seluruh atribut tetap dievaluasi dalam proses pengambilan keputusan, sedangkan metaheuristik hanya menentukan strategi akuisis dan penyimpanan atribut yang paling sesuai.

Evaluasi penelitian dilakukan secara komparatif melalui beberapa konfigurasi uji, yaitu \textit{cloud-only ABAC, edge-fresh ABAC, static cache, adaptive TTL, random search,} dan GWO. Selain itu, \textit{workload} dibagi menjadi skenario rendah, sedang, tinggi, dan sangat dinamis untuk menguji ketahanan mekanisme terhadap perubahan \textit{request rate}, perubahan atribut, \textit{request} kritis, revokasi, dan ketersediaan PIP. Dengan rancangan tersebut, penelitian ini tidak hanya mengukur performa melalui latensi dan \textit{PIP calls}, tetapi juga mengevaluasi kualitas keputusan akses melalui \textit{false allow rate, false deny rate, stale cache hit rate}, dan \textit{revocation false allow rate}.

Dengan dmeikian, Bab II ini menjadi dasar konseptual bagi Bab III, yang akan menjelaskan rancangan arsitektur sistem, model evaluasi ABAC berbasis \textit{edge}, representasi solusi optimasi, fungsi \textit{fitness}, konfigurasi uji, skenario \textit{workload}, serta metrik evaluasi yang digunakan dalam penelitian.

